\section{Código computacional}
\label{ap:codigo}

O notebook completo, contendo a geração das matrizes, as funções de redução de Hessenberg, reflexões de Householder, rotações de Givens, iteração QR, extração dos autovalores, análise de sensibilidade e método da potência, está disponível no repositório:

\begin{center}
  \url{https://github.com/Ederson-Passos/computational-matrix-methods/blob/main/experiments/unit_04_eigenvalue_problems/Computational_Work_4/scripts/problem_1_2.ipynb}
\end{center}

O histórico integral das 818 iterações do método da potência foi exportado para \nolinkurl{iteracoes_metodo_potencia.csv}. Os arquivos \nolinkurl{convergencia_autovalor.pdf}, \nolinkurl{erro_relativo.pdf} e \nolinkurl{matrizes_schur.pdf} contêm figuras geradas a partir dos resultados numéricos.

\subsection{Núcleo do algoritmo de Householder}

\begin{lstlisting}
def house(x):
    x = np.asarray(x, dtype=float).ravel()
    norma = np.linalg.norm(x)
    if np.linalg.norm(x[1:]) == 0:
        return np.zeros((len(x), 1)), 0.0
    sinal = 1.0 if x[0] >= 0 else -1.0
    alpha = -sinal * norma
    v = x.copy()
    v[0] -= alpha
    v = v.reshape(-1, 1)
    c = 2.0 / float(v.T @ v)
    return v, c
\end{lstlisting}

\subsection{Núcleo do método da potência}

\begin{lstlisting}
def metodo_da_potencia(A, x0, tol=1e-8, max_iter=2000):
    q = normalizar_com_sinal(x0.astype(float).copy())
    lam_ant = quociente_rayleigh(q, A)
    historico = []

    for k in range(1, max_iter + 1):
        z = A @ q
        q = normalizar_com_sinal(z)
        lam = quociente_rayleigh(q, A)

        erro = abs(lam - lam_ant) / abs(lam)
        historico.append((k, lam, erro))

        if erro < tol:
            break
        lam_ant = lam

    return lam, q, historico
\end{lstlisting}

O código completo e executável permanece no notebook vinculado acima; os trechos deste apêndice destacam os núcleos dos dois procedimentos numéricos empregados.
